% --------------------------------------------------------------------
% 이론적 배경 섹션입니다.
% Theoretical background section.
%
% 이 파일에서는 핵심 개념, 선행 연구 흐름, 연구 분석틀을 정리합니다.
% This file organizes key concepts, prior research trends, and the analytical framework.
% --------------------------------------------------------------------

% 발표 목차와 사이드바에 표시될 큰 섹션 이름입니다.
% Major section title shown in the table of contents and sidebar.
\section{이론적 배경}

% 핵심 개념 슬라이드의 하위 섹션입니다.
% Subsection for the key-concepts slide.
\subsection{핵심 개념}

% 핵심 개념은 연구 방법에서 실제 분석 기준으로 다시 사용되어야 합니다.
% Key concepts should later be reused as actual criteria in the methods section.
\begin{frame}[t]{핵심 개념 정리}
\small

% 첫 번째 개념 상자입니다. 연구의 중심 개념을 정의합니다.
% First concept box. Define the central concept of the study.
\begin{block}{개념 1: 연구의 중심 개념}
\begin{itemize}
\tightitems

% 정의는 문헌 정의와 본 연구의 조작적 정의를 구분해 쓰면 좋습니다.
% It is helpful to distinguish literature definitions from your operational definition.
\item 개념의 정의와 범위를 간결하게 제시합니다.
\item 본 연구에서 사용하는 조작적 정의를 함께 적습니다.
\end{itemize}
\end{block}

% 두 번째 개념 상자입니다. 분석틀이나 설계 기준을 넣습니다.
% Second concept box. Use it for the analytical framework or design criteria.
\begin{block}{개념 2: 분석 또는 설계의 기준}
\begin{itemize}
\tightitems

% 여기의 기준은 뒤의 연구 문제, 자료, 분석 방법과 연결되어야 합니다.
% These criteria should connect to the later research questions, data, and analysis methods.
\item 연구 문제를 분석할 때 사용할 이론, 기준, 프레임워크를 설명합니다.
\item 필요하면 표나 개념도로 구성 요소를 정리합니다.
\end{itemize}
\end{block}

\slidenote{
이론적 배경은 연구를 장식하는 문헌 목록이 아니라 분석틀을 세우는 부분입니다. 뒤의 연구 방법에서 이 개념들이 어떻게 쓰이는지 연결해 줍니다.
}
\end{frame}

% 선행 연구는 단순 나열보다 연구 흐름과 남은 공백을 중심으로 구성합니다.
% Prior research should emphasize trends and remaining gaps rather than a simple list.
\subsection{선행 연구}
\begin{frame}[t]{선행 연구 흐름}
\small

% 세 개의 열은 "초기 흐름 - 최근 흐름 - 남은 공백"을 비교하기 위한 구조입니다.
% The three columns compare "early trend - recent trend - remaining gap."
\begin{columns}[T]

% 첫 번째 열: 초기 연구 흐름입니다.
% First column: early research trend.
\begin{column}{0.31\textwidth}
\begin{block}{흐름 1}
\scriptsize
초기 연구가 어떤 문제를 제기했는지 정리합니다.
\end{block}
\end{column}

% 두 번째 열: 최근 연구 흐름입니다.
% Second column: recent research trend.
\begin{column}{0.31\textwidth}
\begin{block}{흐름 2}
\scriptsize
최근 연구가 어떤 방법과 결과를 제시했는지 정리합니다.
\end{block}
\end{column}

% 세 번째 열: 본 연구가 다룰 연구 공백입니다.
% Third column: the research gap that this study will address.
\begin{column}{0.31\textwidth}
\begin{alertblock}{남은 공백}
\scriptsize
아직 충분히 다루지 못한 대상, 맥락, 방법, 자료를 제시합니다.
\end{alertblock}
\end{column}
\end{columns}

\vspace{0.25cm}

% 선행 연구의 흐름도, 개념도, 표를 넣기 전까지 사용할 자리표시자입니다.
% Placeholder for a literature map, conceptual diagram, or summary table.
\placeholderbox{0.72\textwidth}{1.6cm}

\slidenote{
선행 연구는 시간순 나열보다 흐름과 공백 중심으로 제시합니다. 마지막에는 반드시 본 연구가 채울 공백을 말합니다.
}
\end{frame}

% 분석틀은 이론적 배경과 연구 방법을 연결하는 다리 역할을 합니다.
% The analytical framework bridges the theoretical background and research methods.
\subsection{연구 분석틀}
\begin{frame}[t]{연구 분석틀}
\small

% 번호 목록은 분석틀의 구성 요소를 순서 있게 설명할 때 적합합니다.
% A numbered list is useful for presenting framework components in order.
\begin{block}{분석틀 구성}
\begin{enumerate}
\tightitems

% 각 항목은 실제 분석에서 어떤 자료와 연결되는지 설명할 수 있어야 합니다.
% Each item should be explainable in terms of the actual data it will analyze.
\item 이론적 기준: 연구 현상을 해석할 개념 또는 프레임워크
\item 자료 기준: 어떤 자료를 증거로 사용할 것인지에 대한 기준
\item 분석 기준: 코딩, 통계, 비교, 검증 방법의 기준
\end{enumerate}
\end{block}

% 강조 상자는 청중이 이 슬라이드에서 기억해야 할 메시지를 요약합니다.
% The highlighted box summarizes what the audience should remember from this slide.
\begin{exampleblock}{발표에서 강조할 점}
분석틀은 연구 문제와 연구 방법 사이의 다리입니다. 이 슬라이드에서 ``무엇을 어떤 관점으로 볼 것인가''를 분명히 설명합니다.
\end{exampleblock}

\slidenote{
분석틀을 설명할 때는 각 범주가 실제 자료 분석에서 어떻게 판별되는지 예시를 들어 주면 좋습니다.
}
\end{frame}
